\section{Átomos positivos y proyección}
Si $U$ prescribe $h$ casillas RC compatibles en la máscara, su probabilidad bajo
las torres es $(\prod_{e\in U}D_e)\per D_{-U}/\per D$. Al reescalar $D_{-U}$
se extraen exactamente los factores de filas y columnas. El factorial, las potencias
de tamaño, las actividades y los factores de escalado se conservan exactamente en
la identidad. Aplicar la parte 04 dos veces cuesta a lo sumo la suma de los errores
en los logaritmos. La cota es independiente de lo pequeña que sea una masa positiva.
Los átomos imposibles quedan fuera de la suma y nunca se logaritman.
\paragraph{Error puro de torres.}
Para la función $G_q$ de \emph{Grados intermedios}, parte 03, se usa además la
identidad de normalización de R9, partes 06--07. Cambiar $\log H_*$ por $\log H_q$ y
truncar Stirling antes del grado $q$ cuesta $(Cq/m)^q+Ce^{-cj}$ en el tamaño físico.
El resto real de Stirling se acota por el primer término omitido con sus coeficientes
de Bernoulli; para $q$ pequeño fijo se cambia $C$. Su polinomio comienza en $u/12$
cuando se incluye ese grado y permanece $O(|u|)$ para $|u|\le c_0/q$.
Recordemos por qué el término constante de $G_q$ desaparece. Con el núcleo balanceado $z$
y $v(1-uv)=z\exp(\alpha+\beta)$, $v$ conserva ambos márgenes uno.
En el centro $v=z$. Diferenciando e integrando se obtiene
$\langle\alpha'+\beta'\rangle=-\int z^2$. Los grados formados con
$\ell=v(1-uv)/(1-2uv)$ tienen derivadas medias que aportan $+\int z^2$
al medio logaritmo de los grados dividido por $u$. Los dos términos se cancelan.
El término local usa $c(y)=\frac12\log(1-2y)-2\log(1-y)=y+O(y^3)$; luego
$(\langle c(uv)\rangle-u)/u^2=O(u)$. En el centro, $W=z-1$ y $\bar\ell=1$,
por lo que $\log d(1)$ cancela el medio logaritmo del determinante del cofactor.
Son identidades bilineales que conservan su valor en los datos balanceados complejos.
Los términos de balanceo, grados y cofactor son holomorfos en un disco fijo por
la parte 02 y la rama de $v$ (para $u$ pequeño su ecuación es una perturbación
analítica de la misma ecuación de márgenes). Como $H_q=d(1)+O(|u|)$ con constante
uniforme, resulta $G_q/u=O(1)$ en un disco $|u|\le c_0/q$, reduciendo $c_0$.
Conocer únicamente el cero en $u=0$ y la cota $|G_q|\le C$ en ese disco daría
un coste $O(q)$ al dividir; la suma explícita de la parte 04 evita esa pérdida.
\paragraph{Qué pasa al activar las marcas.}
En el tamaño físico, para cada configuración de activaciones se tiene un polinomio
en campos libres $\zeta$ con los pesos exactos de sus átomos como coeficientes.
Sustituir cada logaritmo por la aproximación anterior da errores multiplicativos
$O(\Delta_q)$ por átomo, con $\Delta_q=(Cq/j)^q+Ce^{-cj}$.
En el polidisco de campos de \emph{Grados ponderados}, parte 04, el mayorante
absoluto de la suma de pesos aproximados es $e^{Cv}$ y su partición tiene
inversa de módulo $\le e^{Cv}$. Bajo esa entrada distinta, el error relativo
de la suma es a lo sumo $e^{C'v}\Delta_q<1/2$. Así el error de los logaritmos
se prolonga en los campos con la misma cota al tamaño físico.
La diferencia alternada sobre activaciones cuesta como máximo $2^v$ antes de la
proyección en los campos. Sus ganancias adicionales requieren que ambas evaluaciones
coincidan en cada cara de actividad cero. No hemos diferenciado en $u$ el resto físico.
\paragraph{Una comprobación algebraica en una raíz de dos fibras.}
Para dos grupos forestales independientes, $y_i$ es la ocupación binaria de la fibra,
$m_i=\E y_i$, $v_i=m_i(1-m_i)$, $F(y)=y^2-2y$, $L(y)=-y$.
Con $P=F_1F_2$ y $B=P-L_1F_2-F_1L_2$ se tiene $P(y)+B(y)=0$ en cada átomo.
En las medias, $F(m_i)-L(m_i)=-v_i$, luego $(P+B)(m)=-v_1F(m_2)-v_2F(m_1)$.
El término $\frac12\operatorname{tr}(\operatorname{Cov}(y)P''(m))$ aporta el opuesto.
La igualdad persiste al multiplicar por formas lineales en otros grupos distintos
independientes, y al derivar en campos locales que preservan esa independencia.
Esto comprueba una de las raíces excepcionales de \emph{Patrones finitos}, parte 04.
El uso en un bosque completo requiere también la reducción de sus ramas a esas raíces.
